\begin{graphillustration}{Клонирование: $x'$ смежна с $x$ и со всеми соседями $x$ (штриховые рёбра). Справа случай 2: $A$ — класс цвета $x$; множество $(A\setminus\{x\})\cup\{x'\}$ независимо и получает отдельный цвет.}{fig:14-clone}
\begin{tikzpicture}[font=\small]
\begin{scope}
\node[gi vertex] (x) at (0,0) {$x$};
\node[gi vertex] (p) at (-1.2,1.2) {};
\node[gi vertex] (q) at (1.2,1.2) {};
\node[gi vertex] (r) at (0,-1.3) {};
\node[gi selected] (xx) at (1.4,-0.2) {$x'$};
\draw[gi edge] (x) -- (p); \draw[gi edge] (x) -- (q); \draw[gi edge] (x) -- (r); \draw[gi edge] (p) -- (q);
\draw[gi alternate] (xx) -- (x); \draw[gi alternate] (xx) -- (p); \draw[gi alternate] (xx) -- (q); \draw[gi alternate] (xx) -- (r);
\node[gi note] at (0,-2.0) {$G\circ x$};
\end{scope}
\begin{scope}[shift={(5.5,0)}]
\node[gi vertex] (x) at (0,0) {$x$};
\node[gi vertex] (a1) at (-1.6,-1.0) {};
\node[gi vertex] (a2) at (-0.5,-1.4) {};
\node[gi selected] (xx) at (1.4,-0.2) {$x'$};
\node[gi vertex] (n1) at (0,1.2) {};
\draw[gi edge] (x) -- (n1); \draw[gi alternate] (xx) -- (x); \draw[gi alternate] (xx) -- (n1);
\draw[gi edge] (a1) -- (n1); \draw[gi edge] (a2) -- (n1);
\node[gi group,fit=(x)(a1)(a2),label=left:{$A$}] {};
\node[gi note,anchor=west] at (2.2,0.2) {новый цвет:\\$(A\setminus\{x\})\cup\{x'\}$};
\end{scope}
\end{tikzpicture}
\end{graphillustration}
